\chapter{Outputs}
\label{ch:outputs}

\newthought{Everything lands under \file{outputs/<scan>/}.} A run's products are
project-local, layout-stable, and designed to be read by tools you already have:
\textsc{hdf5} and \textsc{csv} for the records, a bucketed artifact tree for per-sample
files, \textsc{json} for the summaries. This chapter is the tour.

\section{The output tree}
\label{sec:output-tree}

\begin{jfiletree}
\dirtree{%
.1 \dtfoldopen{outputs/<scan>/}.
.2 \dtfoldopen{DATABASE/}.
.3 \dtdata{samples.hdf5}\DTcomment{the record of record}.
.3 \dtcsv{samples.csv}\DTcomment{full-column export}.
.2 \dtfoldopen{SAMPLE/}.
.3 \dtfold{000001/}\DTcomment{bucket of $\le$ limit Samples}.
.3 \dtarchive{000002.tar.gz}\DTcomment{packed sealed bucket}.
.2 \dtjson{levelset.json}\DTcomment{AdaptiveBridson runs only}.
.2 \dtjson{run\_summary.json}.
.2 \dtcsv{run\_summary.csv}.
.2 \dtdoc{run\_summary.txt}.
}
\end{jfiletree}

Alongside, at the project root: \file{logs/} (Chapter~\ref{ch:logging}),
\file{images/<scan>/} (the flowchart export and generated figures,
Chapters~\ref{ch:workflow-model} and~\ref{ch:jarvisplot}), and
\file{checkpoints/<scan>/} (Chapter~\ref{ch:checkpoint}).

\section{\code{DATABASE/samples.hdf5}}
\label{sec:database-hdf5}

The primary record: one row per Sample --- succeeded \emph{or failed} --- containing the
full observables dictionary as a \textsc{json} document: parameters, every captured
calculator/opera output, likelihood terms and \code{LogL}, the \textsc{uuid}, status,
timing, and the paths of any \yamlkey{save}d artifacts. Rows are appended in Archiver
batches (\yamlkey{archiver.batch\_size} with a time-based flush), so the file trails the
scan by at most a flush interval.

Reading it directly:

\begin{terminal}
\begin{jcmd}
python3 - <<'PY'
import h5py, json
with h5py.File("outputs/scan-01/DATABASE/samples.hdf5") as f:
    rows = [json.loads(r) for r in f["samples"][...]]
print(len(rows), rows[0]["LogL"])
PY
\end{jcmd}
\end{terminal}

\section{\code{DATABASE/samples.csv}}
\label{sec:database-csv}

A \emph{full-column} flat export of the same records, written at the end of the run:
every observable becomes a column (nested values as \textsc{json} text), so nothing in
the \textsc{hdf5} is missing from the \textsc{csv}. For most analysis,
\code{pandas.read\_csv} on this file is all you need --- and it feeds straight back
into a \sampler{CSV} replay (Chapter~\ref{ch:csv-sampler}).

Missing this file --- an interrupted run, or a deleted \textsc{csv} --- is not a lost
cause: \cli{Jarvis2 convert <task>.yaml} regenerates it (and any other
\file{DATABASE/*.hdf5}) from the \textsc{hdf5} record alone, without touching Redis
(Section~\ref{sec:cli-convert}).

\section{\code{SAMPLE/}: artifacts}
\label{sec:sample-outputs}

Per-sample keepsakes --- \yamlkey{save}d input/output files and the Sample's own log ---
under \file{SAMPLE/<bucket>/<uuid>/}. The bucket policy
(Section~\ref{sec:sample-directory}) bounds directory sizes and packs finished buckets
into \file{<bucket>.tar.gz}; the records in the database always carry the artifact
paths, so locating a specific point's spectrum file is a lookup, not a search.

An opera-only scan that saves nothing keeps \file{SAMPLE/} essentially empty --- the
directories exist only where content does.

\section{\code{run\_summary}: the end-of-run receipt}
\label{sec:run-summary-file}

Written at shutdown in three formats (\file{.json} for machines, \file{.csv} for
spreadsheets, \file{.txt} for humans), with a frozen field order: identity
(\yamlkey{project\_name}, scan, \yamlkey{run\_id}, timestamps), counters (submitted,
completed, failed), and throughput --- \code{wall\_time\_sec}, \code{samples\_per\_sec},
\code{samples\_per\_min}, \code{avg\_sample\_sec}, and per-point evaluation latencies
where Workers reported them. The same numbers appear in the console's
\code{[Scan Performance]} block (Chapter~\ref{ch:monitoring}).

\begin{jnote}
The layout and field order of everything in this chapter is contract-stable
(Section~\ref{sec:contracts}): pipelines that parse \file{samples.csv} or
\file{run\_summary.json} survive upgrades.
\end{jnote}

\begin{jtip}
Archive \file{outputs/<scan>/} and you have archived the scan: records, artifacts,
summary --- plus the card itself if you pack at the project level
(\cli{Jarvis2 project pack . --repro}, Chapter~\ref{ch:project-cli}).
\end{jtip}
